\inicializarSeccion{Interpretación de los resultados inferenciales}

Con los resultados obtenidos de nuestras pruebas estadísticas, podemos determinar si las variaciones observadas en las variables cuantitativas son estadísticamente significativas. Observamos que los datos presentan una variación significativa en el nivel promedio de oxígeno disuelto en el agua, lo cual sugiere que las condiciones del agua están afectando la concentración de oxígeno. Considerando las evaluaciones ideales para el oxigeno disuelto, estos resultados indican que el agua puede no estar en condiciones óptimas para la vida acuática. Hay muchisimos factores que pueden estar influyendo en estos resultados, pero por lo menos en este análisis se observo y se concluyó que la temperatura, la salinidad, la profundidad, y la densidad del agua afecta significativamente la concentración de oxígeno disuelto. El oceano esta en constante cambio y estos factores pueden estar interactuando de manera compleja. La estadistica nos da las herramientas para analizar y comprender estas interacciones de una manera sistemática. Dejandonos así un mejor entendimiento de los procesos que ocurren en el oceano.